\chapter{SQL Level 1: Tables, Rows and Reading Data from One Table}
\companion{Programming with Mosh: SQL Tutorial, Full Database Course for Beginners}{\yts{programming+with+mosh+sql+tutorial+full+database+course+for+beginners}}

\section*{Why this part exists, and how much SQL InfoEdge asks}
The InfoEdge candidate whose experience this book follows said Rounds 1 and 2 are ``drilling'' rounds with three or four questions from each of six
topics, and \textbf{SQL is one of the six} (with Python, machine learning, deep learning, statistics and probability). He did not remember the exact SQL
questions, so this part does not invent ``PYQs''. Instead, questions marked \textcolor{exframe}{\faCheckCircle\ \textbf{Most asked}} are the standard
data-science SQL questions that appear in almost every analytics and data-science interview (Nth highest salary, joins and row counts, WHERE vs HAVING,
ranking functions, duplicates, NULL traps, streaks, funnels, retention). They come first in each chapter's question set.

The part is organised in four levels. Each level assumes only the previous ones.
\begin{center}\small
\begin{tabularx}{\linewidth}{l l X}
\toprule
Level & Chapter & What you can do at the end\\
\midrule
1 & 45 & Read any single table: choose columns, filter rows, sort, handle NULL, create new columns with CASE, strings and dates\\
2 & 46 & Summarise with GROUP BY and HAVING; combine tables with every kind of JOIN; stack results with UNION\\
3 & 47 & Nest queries (subqueries, CTEs, recursive CTEs); use window functions (ranking, LAG/LEAD, running totals); change data safely\\
4 & 48 & Solve the classic interview patterns (Nth highest, top-N per group, duplicates, streaks, funnels, A/B, cohorts, median, pivot);
build leak-free ML feature tables; explain indexes, query plans and schema design\\
\bottomrule
\end{tabularx}
\end{center}

\begin{tcolorbox}[colback=accentlight,colframe=accent,boxrule=0.5pt,arc=2pt]
\small\faLaptopCode\ \textbf{Every query in this part was actually run.} The database is in \code{InfoEdge\_Interview\_Companion/sql\_lab/schema.sql};
the queries are in \code{sql\_lab/queries.sql}; the output printed under each query is the real SQLite output, generated by \code{sql\_lab/run\_queries.py}.
The Colab notebook \code{sql\_lab/SQL\_Lab.ipynb} creates the database inside Python (no installation: SQLite ships with Python) and runs every query, plus
practice exercises with hidden solutions. Type each query yourself before reading its output.
\end{tcolorbox}

\section{What a database is, from zero}
\stuck{What is a relational database? (beginner explanation)}{\yts{what+is+a+relational+database+explained+for+beginners}}
\subsection{Tables, rows and columns}
Think of an Excel workbook. Each sheet holds one kind of thing: one sheet for candidates, one for jobs, one for applications. A \textbf{relational
database} is the same idea made strict and fast:
\begin{itemize}
\item a \textbf{table} is one sheet; it holds one kind of thing (candidates);
\item a \textbf{row} (also \emph{record}) is one thing (one candidate, Aarav);
\item a \textbf{column} (also \emph{field} or \emph{attribute}) is one property that every row has (city, years of experience);
\item every column has a \textbf{data type} that the database enforces: \code{INTEGER} (whole numbers), \code{DECIMAL(5,1)} (numbers with up to 5 digits,
1 after the point), \code{VARCHAR(20)} (text up to 20 characters), \code{DATE} (\code{'2026-03-01'}), \code{BOOLEAN}. Excel would let you type ``five''
into an experience cell; a database refuses.
\end{itemize}
\textbf{SQL} (Structured Query Language, said ``sequel'' or ``S-Q-L'') is the language for asking a database questions. It is \emph{declarative}: you
describe \emph{what} rows you want, not \emph{how} to loop over them; the database's \emph{query optimiser} decides how. That is why one line of SQL can
replace twenty lines of Python loops, and why the same query can run on ten rows or ten billion.

\subsection{NULL: the value that means ``unknown''}
A cell can be empty. SQL marks an empty cell with \code{NULL}, meaning \emph{unknown or not applicable}. NULL is \textbf{not} zero and \textbf{not} the
empty string \code{''}. A fresher has no current salary: NULL. A candidate who skipped the city field: NULL. Because NULL means ``unknown'', any
comparison with NULL is also unknown (Section~\ref{sec:sqlnull}); this single fact causes more wrong interview answers than anything else in SQL.

\subsection{Keys: how rows are identified and linked}
\begin{itemize}
\item A \textbf{primary key} (PK) is a column (or set of columns) whose value is unique for every row and never NULL. \code{cand\_id} identifies a
candidate even if two candidates are both called Aarav.
\item A \textbf{foreign key} (FK) is a column that stores another table's primary key, which creates a link. \code{applications.cand\_id} says
\emph{which} candidate applied. The database can refuse an application for candidate 999 if no candidate 999 exists (\emph{referential integrity}).
\item A \textbf{unique key} is unique like a PK but may allow NULL and a table may have several (e.g.\ email).
\item A \textbf{composite key} uses several columns together: in a skills table, (\code{cand\_id}, \code{skill}) is unique but neither column alone is.
\end{itemize}
Why split data into several linked tables instead of one giant sheet? If the company name were typed into every job and application row, renaming
``PayWave'' would mean editing hundreds of rows and missing one would make the data contradict itself. Storing each fact once and linking by keys is
called \textbf{normalisation} (Chapter~\ref{ch:sql4}).

\section{The practice database: MiniNaukri}
\stuck{SQL primary key and foreign key explained}{\yts{sql+primary+key+foreign+key+explained+simply}}
All queries in Part IX use one small job-portal database, deliberately close to what Naukri stores. It is small enough that you can check every answer by
eye, and it contains the awkward things real data has: NULLs, a duplicate row, a company with no jobs, candidates who never applied, ties in salary.

\begin{center}
\begin{tikzpicture}[tbl/.style={draw=accent,fill=accentlight,rounded corners=2pt,align=left,font=\scriptsize\ttfamily,inner sep=3pt},
  every edge/.style={draw=qframe,-{Latex[length=2mm]}}, font=\scriptsize]
\node[tbl] (co) at (0,0) {\textbf{companies}\\company\_id PK\\name, city\\industry, size};
\node[tbl] (jo) at (4.2,0) {\textbf{jobs}\\job\_id PK\\company\_id FK\\title, city, min\_exp\\max\_ctc\_lpa, posted\_date};
\node[tbl] (ap) at (8.6,0) {\textbf{applications}\\app\_id PK\\cand\_id FK, job\_id FK\\applied\_date, status};
\node[tbl] (ca) at (8.6,-3.1) {\textbf{candidates}\\cand\_id PK, name, city\\exp\_yrs, ctc\_lpa\\signup\_date\\referred\_by FK (self)};
\node[tbl] (lo) at (4.2,-3.1) {\textbf{logins}\\cand\_id FK\\login\_date};
\node[tbl] (im) at (12.6,-3.1) {\textbf{impressions}\\imp\_id PK, cand\_id\\job\_id, shown\_date\\variant, clicked};
\node[tbl] (em) at (0,-3.1) {\textbf{employees}\\emp\_id PK, name, dept\\salary\\manager\_id FK (self)\\hire\_date};
\draw[draw=qframe,-{Latex[length=2mm]}] (jo) -- node[above]{many:1} (co);
\draw[draw=qframe,-{Latex[length=2mm]}] (ap) -- node[above]{many:1} (jo);
\draw[draw=qframe,-{Latex[length=2mm]}] (ap) -- node[right]{many:1} (ca);
\draw[draw=qframe,-{Latex[length=2mm]}] (lo) -- node[above]{many:1} (ca);
\end{tikzpicture}
\end{center}
An arrow points from the table holding the foreign key to the table it refers to. ``many:1'' means many jobs can belong to one company, but each job
belongs to exactly one company.

\begin{center}\small
\begin{tabularx}{\linewidth}{l r X}
\toprule
Table & Rows & What one row is, and the awkward bits planted in it\\
\midrule
companies & 6 & A hiring company. GreenGrid (Delhi) has posted no jobs.\\
candidates & 10 & A job seeker. Sana's city is NULL; Kabir is a fresher with NULL salary; Sana and Priya never applied; \code{referred\_by} points to
another candidate.\\
jobs & 10 & A job posting. Job 210 (Hyderabad) received no applications.\\
applications & 23 & One candidate applying to one job, with the current status. Row 22 is an accidental duplicate of row 1.\\
logins & 18 & One login event. Candidate 109 logged in twice on 4 March.\\
impressions & 16 & A job card shown to a candidate during an A/B test of a new ranking (variant B), and whether it was clicked.\\
employees & 11 & An internal employee, for the classic salary questions. Zoya's salary is NULL; Kavya and Manav tie at 45.\\
\bottomrule
\end{tabularx}
\end{center}

\subsection{Running SQL from Python (the way you will in a notebook round)}
\begin{lstlisting}
import sqlite3, pandas as pd
con = sqlite3.connect(":memory:")                 # a database that lives in RAM
con.executescript(open("sql_lab/schema.sql").read())   # create and fill the tables
df = pd.read_sql("SELECT * FROM candidates WHERE exp_yrs >= 5", con)
print(df)
\end{lstlisting}
\code{pd.read\_sql} returns a DataFrame, so you can move between SQL and pandas freely. In industry the connection is to MySQL, PostgreSQL, Redshift,
BigQuery or Hive rather than SQLite, but the SQL in this part is standard and runs on all of them (dialect differences are listed in
Section~\ref{sec:dialects}).

\section{SELECT and FROM: choosing columns}
\stuck{SQL SELECT statement for beginners}{\yts{sql+select+statement+tutorial+for+beginners}}
The smallest query names a table (\code{FROM}) and the columns to show (\code{SELECT}). \code{*} means ``all columns''.
\runsql{L1_select_all}
\runsql{L1_columns}
Read the output like a small spreadsheet: one line per row, columns separated by \code{|}. Note \code{NULL} in Sana's city: the value is unknown.

\begin{trap}[: SELECT * in production]
\code{SELECT *} is fine for exploring. In production code and in interviews, list the columns you need: it transfers less data, documents intent, and
does not break silently when someone adds a column. On columnar warehouses (BigQuery, Redshift) you also pay per column read.
\end{trap}

\section{WHERE: choosing rows}
\stuck{SQL WHERE clause, AND OR NOT, IN, BETWEEN, LIKE}{\yts{sql+where+clause+and+or+in+between+like+tutorial}}
\code{WHERE} keeps only the rows for which a condition is true. Comparison operators: \code{=}, \code{<>} (or \code{!=}, ``not equal''), \code{<},
\code{<=}, \code{>}, \code{>=}. Text values go in single quotes.
\runsql{L1_where}

\subsection{AND, OR and the precedence trap}
\code{AND} needs both conditions true; \code{OR} needs at least one. \textbf{AND is evaluated before OR}, exactly as multiplication comes before addition.
Suppose the recruiter asks for ``candidates from Bengaluru or Mumbai with at least 5 years''. The obvious query is wrong:
\runsql{L1_and_or}
SQL read it as \code{city = 'Bengaluru' OR (city = 'Mumbai' AND exp\_yrs >= 5)}, so every Bengaluru candidate came back, even with 2 years (Aarav,
Ishaan). Brackets fix it:
\runsql{L1_and_or_fixed}
\begin{recap}
Whenever a WHERE mixes AND and OR, add brackets, even if you think precedence already does what you want. Interviewers plant this bug in ``find the error''
questions.
\end{recap}

\subsection{IN, BETWEEN and LIKE}
\begin{itemize}
\item \code{x IN ('a','b','c')} is a short way to write \code{x = 'a' OR x = 'b' OR x = 'c'}.
\item \code{x BETWEEN 3 AND 10} means \code{x >= 3 AND x <= 10}: \textbf{both ends included}.
\item \code{x LIKE 'Data\%'} matches text patterns: \code{\%} = any number of characters (including none), \code{\_} = exactly one character.
\code{'\%Engineer'} = ends with Engineer; \code{'\%data\%'} = contains data; \code{'\_\_\_'} = exactly three characters.
\end{itemize}
\runsql{L1_in_between}
\runsql{L1_like}

\begin{trap}[: BETWEEN with timestamps, and LIKE's case sensitivity]
\begin{itemize}
\item If \code{applied\_at} is a \emph{timestamp}, \code{applied\_at BETWEEN '2026-03-01' AND '2026-03-31'} stops at 31 March 00:00:00 and silently drops
everything on the afternoon of the 31st. Use a half-open range: \code{applied\_at >= '2026-03-01' AND applied\_at < '2026-04-01'}.
\item \code{LIKE} is case-insensitive for English letters in SQLite and (by default collation) MySQL, but \textbf{case-sensitive in PostgreSQL}
(\code{ILIKE} is the insensitive version there). The portable way: \code{LOWER(title) LIKE '\%data\%'}.
\end{itemize}
\end{trap}

\section{NULL and three-valued logic}\label{sec:sqlnull}
\stuck{SQL NULL values explained, IS NULL and three-valued logic}{\yts{sql+null+three+valued+logic+explained}}
Because NULL means ``unknown'', SQL logic has three values: TRUE, FALSE and UNKNOWN. Is an unknown city equal to Delhi? Unknown. Is an unknown city
\emph{not} equal to Delhi? Also unknown. Is NULL equal to NULL? Two unknowns might or might not be the same, so: unknown. \textbf{WHERE keeps a row only
when the condition is TRUE}; UNKNOWN rows are dropped, exactly like FALSE rows.

\begin{center}\small
\begin{tabular}{c c c c}
\toprule
$p$ & $q$ & $p$ AND $q$ & $p$ OR $q$\\
\midrule
TRUE & UNKNOWN & UNKNOWN & TRUE\\
FALSE & UNKNOWN & FALSE & UNKNOWN\\
UNKNOWN & UNKNOWN & UNKNOWN & UNKNOWN\\
\bottomrule
\end{tabular}
\qquad
\begin{tabular}{l c}
\toprule
Expression & Result\\
\midrule
\code{NULL = NULL} & UNKNOWN\\
\code{NULL <> 5} & UNKNOWN\\
\code{NULL + 5} & NULL\\
\code{NOT UNKNOWN} & UNKNOWN\\
\code{NULL IS NULL} & TRUE\\
\bottomrule
\end{tabular}
\end{center}
Rule of thumb: FALSE AND anything is FALSE; TRUE OR anything is TRUE; otherwise UNKNOWN spreads.

So \code{= NULL} never finds anything. You must use \code{IS NULL} / \code{IS NOT NULL}:
\runsql{L1_null_wrong}
\runsql{L1_null_right}
The more dangerous version hides inside ordinary filters. There are 10 candidates and 3 live in Bengaluru. How many are ``not in Bengaluru''?
\runsql{L1_not_equal_null}
Six, not seven: Sana's NULL city makes \code{city <> 'Bengaluru'} UNKNOWN, so she is dropped. If the business question is ``everyone except Bengaluru'',
write \code{WHERE city <> 'Bengaluru' OR city IS NULL}, or \code{WHERE COALESCE(city, '') <> 'Bengaluru'}.

\subsection{COALESCE: replacing NULL}
\code{COALESCE(a, b, c)} returns the first argument that is not NULL. It is SQL's \code{fillna}.
\runsql{L1_coalesce}
\begin{practical}[: is ``fill with 0'' right?]
Filling Kabir's salary with 0 is correct for ``total payroll'' but wrong for ``average salary'' (it pulls the average down) and wrong as an ML feature
unless you also add a flag \code{ctc\_missing} (Chapter~\ref{ch:sql4} builds exactly that). Missingness here means ``fresher'', which is information,
the same lesson as the EDA chapters' informative missingness.
\end{practical}

\section{ORDER BY, LIMIT and OFFSET: sorting and taking the top}
\stuck{SQL ORDER BY and LIMIT tutorial}{\yts{sql+order+by+limit+offset+tutorial}}
\code{ORDER BY col} sorts ascending (\code{ASC}, the default); \code{DESC} sorts descending. You can sort by several columns:
\code{ORDER BY title, city} sorts by title and breaks ties by city. \code{LIMIT n} keeps the first $n$ rows; \code{OFFSET k} skips $k$ rows first (used
for pages: page 2 of size 3 is \code{LIMIT 3 OFFSET 3}).
\runsql{L1_order_limit}
\runsql{L1_offset}
Where does NULL go when sorting? It depends on the database:
\runsql{L1_order_nulls}
SQLite and MySQL treat NULL as the smallest value (first in ASC, last in DESC). PostgreSQL and Oracle treat it as the largest (last in ASC). Standard SQL
lets you say exactly what you want: \code{ORDER BY ctc\_lpa ASC NULLS LAST}.

\begin{trap}[: rows have no order unless you ask]
A table is a \emph{set} of rows: it has no built-in order. Without ORDER BY, \code{LIMIT 3} returns \emph{some} 3 rows, which may change from run to run.
``Top 3'' always needs ORDER BY, and a tie-breaker (\code{ORDER BY ctc\_lpa DESC, cand\_id}) if ties are possible.
\end{trap}

\section{DISTINCT: unique values}
\code{DISTINCT} removes duplicate result rows. With several columns it removes duplicate \emph{combinations}.
\runsql{L1_distinct}
\runsql{L1_distinct_pair}
Three distinct titles, but nine distinct (title, city) pairs. \code{DISTINCT} applies to the whole row, never to one column only.

\section{Computed columns, aliases and integer division}
\stuck{SQL alias and calculated columns}{\yts{sql+calculated+columns+alias+as+tutorial}}
SELECT can compute new columns with arithmetic (\code{+ - * /}) and name them with \code{AS} (an \emph{alias}).
\runsql{L1_computed}
Meera earns 30 LPA, i.e.\ $30 \times 100000 / 12 = 250000$ rupees a month; a 30\% hike would be 39 LPA. Now a trap that costs candidates points:
\runsql{L1_intdiv}
In SQLite, PostgreSQL and SQL Server, \textbf{integer divided by integer is integer division} ($7/2 = 3$). A conversion rate computed as
\code{offers / applications} with two integer counts returns 0 for every row with fewer offers than applications. Multiply by \code{1.0} (or
\code{CAST(x AS DECIMAL)}) first. MySQL is the exception: its \code{/} always returns a decimal.

\section{CASE WHEN: if--else inside a query}
\stuck{SQL CASE WHEN statement explained}{\yts{sql+case+when+statement+explained}}
\code{CASE} checks conditions top to bottom and returns the value of the first one that is true; \code{ELSE} catches everything else (without ELSE you get
NULL). It is how you bin a numeric column, exactly like \code{pd.cut}.
\runsql{L1_case}
Because the conditions are checked in order, the second branch \code{exp\_yrs <= 3} does not need \code{exp\_yrs > 0}: zero was already caught.

\section{String and date functions}
\stuck{SQL string functions and date functions tutorial}{\yts{sql+string+functions+date+functions+tutorial}}
\runsql{L1_strings}
\code{||} joins strings (\code{CONCAT(a, b)} in MySQL). \code{SUBSTR(s, start, length)} counts from 1, not 0.
\runsql{L1_dates}
\code{strftime('\%Y-\%m', d)} extracts year and month (the grouping key for monthly reports); \code{julianday} turns a date into a day number so two dates
can be subtracted; \code{date(d, '+30 days')} adds days.

\subsection{Dialect cheat-sheet}\label{sec:dialects}
The core of SQL is standard. Dates, string joining and a few limits differ. Say in the interview which dialect you are writing; the interviewer will not
mind the small differences.
\begin{center}\scriptsize
\begin{tabularx}{\linewidth}{l X X X}
\toprule
Task & SQLite (this book) & MySQL 8 & PostgreSQL\\
\midrule
Year-month & \code{strftime('\%Y-\%m', d)} & \code{DATE\_FORMAT(d, '\%Y-\%m')} & \code{TO\_CHAR(d, 'YYYY-MM')} or \code{DATE\_TRUNC('month', d)}\\
Days between & \code{julianday(a) - julianday(b)} & \code{DATEDIFF(a, b)} & \code{a - b} (dates)\\
Add 7 days & \code{date(d, '+7 days')} & \code{DATE\_ADD(d, INTERVAL 7 DAY)} & \code{d + INTERVAL '7 days'}\\
Join strings & \code{a || b} & \code{CONCAT(a, b)} & \code{a || b} or \code{CONCAT}\\
First $n$ rows & \code{LIMIT n} & \code{LIMIT n} & \code{LIMIT n} (SQL Server: \code{TOP n})\\
7/2 & 3 & 3.5000 & 3\\
NULL in ASC sort & first & first & last\\
Case-insensitive LIKE & default & default & \code{ILIKE}\\
\bottomrule
\end{tabularx}
\end{center}

\section{The logical order of a query}\label{sec:sqlorder}
\stuck{SQL order of execution explained}{\yts{sql+query+order+of+execution+explained}}
We \emph{write} a query in the order SELECT, FROM, WHERE, GROUP BY, HAVING, ORDER BY, LIMIT. The database \emph{evaluates} it (logically) in a different
order. Knowing this order answers half of all ``why does this error?'' questions.
\begin{center}
\begin{tikzpicture}[st/.style={draw=accent,fill=accentlight,rounded corners=2pt,font=\scriptsize\bfseries,minimum height=7mm,inner sep=3pt},
  node distance=2.5mm]
\node[st] (a) {1 FROM / JOIN};
\node[st,right=of a] (b) {2 WHERE};
\node[st,right=of b] (c) {3 GROUP BY};
\node[st,right=of c] (d) {4 HAVING};
\node[st,right=of d] (e) {5 SELECT (window fns)};
\node[st,right=of e] (f) {6 DISTINCT};
\node[st,below=4mm of f] (g) {7 ORDER BY};
\node[st,left=of g] (h) {8 LIMIT / OFFSET};
\foreach \x/\y in {a/b,b/c,c/d,d/e,e/f,f/g,g/h} \draw[-{Latex[length=1.5mm]}] (\x) -- (\y);
\end{tikzpicture}
\end{center}
Consequences you can now explain from first principles:
\begin{itemize}
\item An alias defined in SELECT (step 5) cannot be used in WHERE (step 2), because it does not exist yet. It \emph{can} be used in ORDER BY (step 7).
(SQLite and MySQL are lenient and sometimes accept it; PostgreSQL and SQL Server reject it. Do not rely on it.)
\item Aggregates like \code{COUNT(*)} cannot appear in WHERE (rows have not been grouped yet); filters on aggregates go in HAVING (step 4).
\item Window functions are computed in step 5, so they cannot be used in WHERE either; wrap the query and filter outside (Chapter 47).
\item ORDER BY can sort by a column you did not select (except after DISTINCT), and LIMIT is applied last.
\end{itemize}

\begin{example}[: walk one query through the eight steps]
\code{SELECT city, COUNT(*) AS n FROM candidates WHERE exp\_yrs >= 2 GROUP BY city HAVING COUNT(*) >= 2 ORDER BY n DESC;}
\begin{enumerate}
\item FROM: all 10 candidates.
\item WHERE \code{exp\_yrs >= 2}: Kabir (0) and Sana (1) removed; 8 left.
\item GROUP BY city: Bengaluru \{Aarav, Meera, Ishaan\}, Mumbai \{Diya, Ananya\}, Delhi \{Priya\}, Pune \{Rohan\}, Noida \{Vikram\}.
\item HAVING \code{COUNT(*) >= 2}: Bengaluru (3) and Mumbai (2) survive.
\item SELECT: (Bengaluru, 3), (Mumbai, 2).
\item DISTINCT: not used, nothing happens.
\item ORDER BY n DESC: Bengaluru first, then Mumbai. (8. no LIMIT.)
\end{enumerate}
\end{example}

% ======================================================================
\section{Interview questions: Level 1}
\stuck{Top SQL interview questions for data analysts and data scientists}{\yts{sql+interview+questions+for+data+scientist+data+analyst}}

\begin{qa}{\classic What is the order in which a SQL query is executed, and why can't I use a SELECT alias in WHERE?}
\oneline{Logically: FROM/JOIN, WHERE, GROUP BY, HAVING, SELECT, DISTINCT, ORDER BY, LIMIT; aliases are created in SELECT, which runs after WHERE, so WHERE
cannot see them.}
\why The written order (SELECT first) is for humans; the logical order is how the result is defined. FROM builds the set of rows (joining tables), WHERE
removes rows, GROUP BY collapses the survivors into groups, HAVING removes groups, SELECT computes the output columns (including aliases and window
functions), DISTINCT removes duplicate output rows, ORDER BY sorts, LIMIT cuts. Each step only sees what earlier steps produced. ``Logical'' matters:
the optimiser may physically push filters earlier or use an index, but the \emph{result} must be as if this order were followed.
\worked \code{SELECT ctc\_lpa*1.3 AS expected FROM candidates WHERE expected > 20} fails in PostgreSQL (``column expected does not exist''). Fix by
repeating the expression (\code{WHERE ctc\_lpa*1.3 > 20}) or by wrapping: \code{SELECT * FROM (SELECT ctc\_lpa*1.3 AS expected FROM candidates) t
WHERE expected > 20}. In ORDER BY the alias works: \code{ORDER BY expected DESC}.
\pushes \emph{Then why does my alias work in WHERE on MySQL?} MySQL and SQLite add a non-standard convenience; portable SQL must not rely on it.
\emph{Where do window functions fit?} Step 5, which is why \code{WHERE ROW\_NUMBER() OVER(...) = 1} is illegal.
\pitfall Saying ``SQL executes top to bottom as written''.
\end{qa}

\begin{qa}{\classic Why does \code{WHERE city = NULL} return nothing, and how many rows does \code{WHERE city <> 'Bengaluru'} return in MiniNaukri?}
\oneline{Any comparison with NULL is UNKNOWN, and WHERE keeps only TRUE rows; so \code{= NULL} matches nothing, and \code{<> 'Bengaluru'} returns 6, not
7, because Sana's NULL city is silently dropped.}
\why NULL means ``value unknown''. Asking ``is unknown equal to unknown?'' cannot be answered, so SQL returns UNKNOWN (third truth value). WHERE treats
UNKNOWN as ``not true'' and discards the row. \code{IS NULL} is a different operator that asks ``is this value missing?'', which has a definite answer.
\worked 10 candidates, 3 in Bengaluru, 1 with NULL city: \code{<> 'Bengaluru'} gives $10 - 3 - 1 = 6$ (query \code{L1\_not\_equal\_null} above).
To include the unknown city: \code{WHERE city <> 'Bengaluru' OR city IS NULL} gives 7.
\pushes \emph{What does \code{COUNT(city)} return?} 9: aggregate functions skip NULLs, only \code{COUNT(*)} counts rows. \emph{What is
\code{NULL + 5}?} NULL; arithmetic with unknown is unknown. \emph{Is there a NULL-safe equality?} Yes: \code{IS NOT DISTINCT FROM} (standard,
PostgreSQL), \code{<=>} (MySQL), \code{IS} (SQLite).
\pitfall Using \code{= NULL}, or forgetting that \code{NOT IN} with a NULL in the list returns nothing (Chapter 47).
\end{qa}

\begin{qa}{\classic Primary key vs unique key vs foreign key: what is the difference?}
\oneline{A primary key uniquely identifies each row and cannot be NULL (one per table); a unique key enforces uniqueness but a table can have several and
they may allow NULL; a foreign key references another table's key and enforces that the referenced row exists.}
\why PK: the identity of the row; the database usually builds an index on it automatically, and other tables point at it. Unique key: a business rule
(``no two candidates share an email''); several are allowed. FK: a link; it prevents orphans (an application for a job that does not exist) and can
cascade deletes (\code{ON DELETE CASCADE}).
\worked \code{candidates.cand\_id} is the PK; \code{applications.cand\_id} is an FK to it; \code{candidates.referred\_by} is an FK to the
\emph{same} table (a self-reference), which is how referral chains are stored. In a skills table, (\code{cand\_id}, \code{skill}) is a composite PK.
\pushes \emph{Surrogate vs natural key?} A surrogate key is a meaningless id (\code{cand\_id = 101}); a natural key is real data (PAN number). Surrogates
never change, natural keys sometimes do (people change phone numbers), so surrogates are preferred as PKs with a unique constraint on the natural key.
\pitfall Saying a unique key ``is the same as a primary key''.
\end{qa}

\begin{qa}{\classic DELETE vs TRUNCATE vs DROP?}
\oneline{DELETE removes chosen rows (with WHERE) and can be rolled back; TRUNCATE removes all rows quickly and keeps the table; DROP removes the table
itself, structure included.}
\why DELETE is a row-by-row data operation (DML): it fires triggers, logs each row, honours WHERE. TRUNCATE is a structural operation (DDL in most
databases): it deallocates the data pages, so it is much faster on big tables, resets auto-increment counters, has no WHERE, and in MySQL and Oracle
cannot be rolled back (PostgreSQL can roll it back inside a transaction). DROP deletes data, indexes, constraints and the definition.
\worked Removing the duplicate application: \code{DELETE FROM applications WHERE app\_id = 22}. Clearing a staging table before each nightly load:
\code{TRUNCATE TABLE staging\_apps}. Retiring an old table: \code{DROP TABLE old\_apps}.
\pushes \emph{Which is fastest?} DROP $\approx$ TRUNCATE $\gg$ DELETE on large tables. \emph{Can TRUNCATE work if other tables reference this one by
FK?} Usually not without disabling or cascading the constraint.
\pitfall ``TRUNCATE is just DELETE without WHERE'' (it is implemented very differently).
\end{qa}

\begin{qa}{Write a query: candidates from Bengaluru or Mumbai with at least 5 years of experience. Then explain what goes wrong without brackets.}
\oneline{\code{WHERE (city = 'Bengaluru' OR city = 'Mumbai') AND exp\_yrs >= 5}, or \code{WHERE city IN ('Bengaluru','Mumbai') AND exp\_yrs >= 5};
without brackets AND binds first and every Bengaluru candidate is returned.}
\why AND has higher precedence than OR, as $\times$ over $+$. \code{A OR B AND C} means \code{A OR (B AND C)}.
\worked The unbracketed query returned Aarav, Diya, Meera, Ishaan; the correct one returns Diya and Meera (outputs \code{L1\_and\_or} and
\code{L1\_and\_or\_fixed}).
\pushes \emph{Why prefer IN?} Shorter, no precedence risk, and easy to extend. \emph{What about NOT?} NOT binds tighter than AND; \code{NOT a AND b} is
\code{(NOT a) AND b}.
\pitfall Testing on data where the bug happens not to show (if no Bengaluru candidate had fewer than 5 years you would never notice).
\end{qa}

\begin{qa}{CHAR vs VARCHAR, and which numeric type should store salary or a rate?}
\oneline{CHAR(n) is fixed length (padded with spaces), VARCHAR(n) variable length; money and rates should be DECIMAL/NUMERIC (exact), not FLOAT
(approximate).}
\why CHAR suits values that always have the same length (country code \code{'IN'}, variant \code{'A'}); VARCHAR saves space for names and titles. FLOAT
stores binary approximations: $0.1 + 0.2 = 0.30000000000000004$, so summing millions of salaries in FLOAT drifts and equality tests fail. DECIMAL(p,s)
stores exact decimal digits. Use INTEGER for counts and ids, DATE/TIMESTAMP for time (never text, which sorts wrongly unless in ISO format), BOOLEAN for
flags.
\worked \code{ctc\_lpa DECIMAL(5,1)} holds up to 9999.9 lakh with one decimal; \code{variant CHAR(1)}.
\pushes \emph{Why does ISO format ('2026-03-01') matter if dates are text?} Because then alphabetical order equals time order; '01/03/2026' sorts wrongly.
\pitfall Storing dates as \code{'1/3/26'} strings.
\end{qa}

\begin{qa}{Find all job titles that contain the word ``data'' in any case, and say how your answer changes across databases.}
\oneline{\code{WHERE LOWER(title) LIKE '\%data\%'} is portable; PostgreSQL also offers \code{ILIKE}.}
\why \code{\%} matches any run of characters. LIKE's case sensitivity depends on the database and its collation: SQLite and MySQL (default) ignore case,
PostgreSQL does not. Lower-casing both sides removes the ambiguity.
\worked In MiniNaukri, 7 of 10 jobs match (all Data Scientist and Data Analyst jobs, output \code{L1\_like}). A leading \code{\%} also prevents
the database from using an ordinary index on \code{title} (Chapter~\ref{ch:sql4}): fine for 10 rows, slow for 50 million.
\pushes \emph{How would you search job descriptions at Naukri scale?} With a full-text index (MySQL FULLTEXT, PostgreSQL \code{tsvector}) or a search
engine such as Elasticsearch, not LIKE.
\pitfall Forgetting that \code{\_} is a wildcard: \code{LIKE 'ML\_\%'} matches ``MLE'' too; escape it when you mean a literal underscore.
\end{qa}

\begin{qa}{\deep SQL is called ``declarative''. What does that mean, and why does it matter for a data scientist?}
\oneline{You state \emph{what} result you want and the database chooses \emph{how} to compute it, so the same query can be run by very different physical
plans, and performance depends on the optimiser, indexes and data layout rather than on how you ordered your loops.}
\why In pandas you write a sequence of operations and they run in that order. In SQL you describe the result set: ``rows from these tables matching
this condition, grouped this way''. The query optimiser estimates the cost of alternative plans (which table to read first, whether to use an index,
nested-loop vs hash join) using table statistics and picks the cheapest. Consequences: (1) two queries that mean the same thing usually run equally fast;
(2) you speed queries up by giving the optimiser options (indexes, sargable filters, fewer columns) rather than by reordering clauses; (3) you inspect
the chosen plan with \code{EXPLAIN} (Chapter~\ref{ch:sql4}).
\worked \code{SELECT * FROM applications WHERE cand\_id = 105}: without an index the plan is \code{SCAN applications} (read all rows); after
\code{CREATE INDEX} it is \code{SEARCH ... USING INDEX} (jump to the matching rows). The query text did not change.
\pushes \emph{So is the logical order irrelevant?} No: it defines the meaning (what is legal, what each clause sees); the physical order is the
optimiser's business. \emph{When would you still push work into SQL instead of pandas?} When data is large: filtering and aggregating in the database
moves megabytes instead of gigabytes over the network.
\pitfall Claiming that putting the most selective condition first in WHERE makes the query faster (optimisers reorder predicates).
\end{qa}

\begin{qa}{\deep A recruiter-analytics dashboard shows the average CTC of Mumbai candidates as 16.5 LPA. A colleague's Python script says 11 LPA for the
same candidates. Both read the same table. What could explain it?}
\oneline{Most likely NULL handling: SQL's AVG ignores NULLs, while the script filled missing salaries with 0 (or the reverse); other suspects are integer
division, duplicates, or different filters on NULL cities.}
\why AVG computes sum of non-NULL values / number of non-NULL values. \code{df['ctc'].fillna(0).mean()} divides by all rows. If 1 of 3 Mumbai candidates
had NULL salary with the others at 18 and 15, SQL gives $33/2 = 16.5$; the script gives $33/3 = 11$. Neither is ``wrong''; they answer different
questions (average among those who reported vs average with unknown treated as zero, which is rarely sensible).
\worked The \code{employees} table shows the same effect: \code{AVG(salary)} = 53 (10 known salaries), while \code{SUM(salary)/COUNT(*)} = 48.18 (11
rows, Zoya unknown) (output \code{L2\_agg\_basic} in Chapter 46).
\pushes \emph{How do you settle it?} Decide the definition with the business, write it down (``average over candidates with a reported CTC''), and report
the coverage alongside (``2 of 3 candidates reported''). \emph{Other causes?} A join that duplicated rows (fan-out, Chapter 46), or \code{city <> ...}
filters that silently dropped NULL cities.
\pitfall Picking one number without explaining the definitional difference.
\end{qa}

\begin{qa}{What is the difference between \code{COUNT(*)}, \code{COUNT(col)} and \code{COUNT(DISTINCT col)}? Give the values for MiniNaukri.}
\oneline{\code{COUNT(*)} counts rows; \code{COUNT(col)} counts rows where col is not NULL; \code{COUNT(DISTINCT col)} counts different non-NULL values.}
\why Aggregates skip NULLs, except \code{COUNT(*)}, which does not look at any column. DISTINCT inside the function deduplicates before counting.
\worked \code{employees}: \code{COUNT(*)} = 11, \code{COUNT(salary)} = 10 (Zoya NULL). \code{applications}: \code{COUNT(*)} = 23 applications,
\code{COUNT(DISTINCT cand\_id)} = 8 candidates, \code{COUNT(DISTINCT job\_id)} = 9 jobs (output \code{L2\_count\_distinct} in Chapter 46).
\pushes \emph{Is \code{COUNT(1)} faster than \code{COUNT(*)}?} No; every modern optimiser treats them identically. \emph{Count distinct pairs?}
\code{COUNT(DISTINCT cand\_id || '-' || job\_id)} or count rows of a \code{SELECT DISTINCT} subquery.
\pitfall Using \code{COUNT(col)} to count rows when col has NULLs.
\end{qa}
